\section{Introduction}\label{sec:introduction}

Weisfeiler--Leman refinement compares graphs by repeatedly refining colors
of ordered vertex tuples. For every fixed dimension, this procedure is
computable in polynomial time. When the dimension is part of the input,
however, the familiar \(n^{O(k)}\) tuple computation is no longer a
polynomial-time algorithm. We determine the complexity of deciding its
answer, both on general graphs and on connected graphs of maximum degree
three.

We use the \emph{joint-update} convention. Initially a \(k\)-tuple records
the equalities and adjacencies among its entries. For each possible
replacement vertex \(z\), form the vector of \(k\) old colors obtained by
replacing, in turn, each coordinate by that same \(z\). The new color
records the old color and the multiset of these entire vectors. Color
names are common to the two graphs. Write \(G\equiv_k H\) if their
\(k\)-tuple color histograms agree at every round.
\Cref{def:joint-wl} gives the recursive definition.

The language \(\WL\) consists of the encodings \((G,H,k)\) with
\(G\equiv_k H\), where \(G,H\) are finite simple uncolored undirected
graphs represented by explicit adjacency matrices and \(k\ge2\) is
encoded in binary. Its restricted version \(\SubWL\) additionally
requires the two graphs to be connected, of the same positive order,
and of maximum degree at most three. Malformed encodings or failures
of the specified graph and dimension conditions are excluded.

\begin{theorem}\label{thm:main}
Both \(\WL\) and \(\SubWL\) are \(\EXPTIME\)-complete under deterministic
polynomial-time many-one reductions.
\end{theorem}

The degree condition is an upper bound: the constructed subcubic graphs
are not required to be regular. The two problems concern equivalence
of a supplied pair at the supplied dimension, rather than finding the
least dimension that identifies one graph among all possible mates.
The companion \citep[Theorem~1.1]{CompanionIdentification}
proves \(\EXPTIME\)-completeness for deciding whether a supplied
dimension identifies one graph among all mates, using our uniform
computation interface together with a separate all-mates reconstruction.
The dimension may exceed the graph order; the upper bound proves that
it can be capped before any tuple table is allocated.

\subsection{History and the role of the degree restriction}

The refinement method originates in Weisfeiler and Leman's work on
graph canonization and its associated algebra \citep{WeisfeilerLeman1968}.
Cai, F{\"u}rer and Immerman established linear lower bounds on the
number of variables required for graph identification. Their examples
include nonisomorphic degree-three graphs with color classes of size at most
four that remain indistinguishable with linearly many counting variables
\citep[Theorem~5.2 and Corollary~6.5]{CaiFurerImmerman1992}.
These are dimension lower bounds, not lower bounds on every algorithm
deciding the outcome of refinement. The joint-update convention and its
bijective \((k+1)\)-pair game characterization are recorded explicitly by
Kiefer and Neuen \citep[Section~2.2 and Theorem~1]{KieferNeuen2019}.
We include the complete histogram-to-empty-game proof because the
extra pair and the timing of the announced bijection are essential.

For dimension supplied as input, Seppelt proved coNP-hardness of
pair equivalence and recorded Berkholz's question of
\(\EXPTIME\)-completeness
\citep[Theorem~5 and the preceding question, p.~82:3]{Seppelt2024}.
Lichter, Ra{\ss}mann and Schweitzer independently obtained coNP-hardness
for simple uncolored graphs
\citep[Theorem~19]{LichterRassmannSchweitzer2025}.
Their work also studies the distinct identification-dimension problem.
\Cref{thm:main} answers the equivalence question for the joint-update
convention and, in addition, for the connected subcubic restriction.

The latter restriction is instructive because Luks proved that
isomorphism of graphs of every fixed maximum degree is decidable in
polynomial time \citep{Luks1982}. Refinement at a prescribed dimension
asks a different question: nonisomorphic graphs can still be equivalent.
Our degree-three result concerns that query, not bounded-degree
graph-isomorphism hardness.

Grohe, Lichter, Neuen and Schweitzer compress CFI constructions into
smaller graphs while preserving long refinement sequences
\citep{cfiCompression2025}. The product-address,
scalar-image, and corner-potential methods are shared with the companion
\emph{Unconditional time lower bounds for Weisfeiler--Leman equivalence}
\citep[Sections~3--4]{CompanionFixedDimension}. That paper fixes the
dimension and proves a quantitative lower bound for arbitrary
deterministic deciders. Here the number of binary coordinates grows with
the source input. We therefore prove the required computation and
projection statements locally and account for every dependence on that
number. A fixed-dimension theorem with dimension-dependent constants
would not supply a uniform polynomial reduction.

\subsection{How the reduction uses the bijective game}

The game formulation on equal-order graphs explains what the reduction
must preserve.
With \(k+1\) named pair slots, each occupied slot holds one vertex
from each graph. Spoiler may discard pairs and request a vacant slot.
Duplicator must then announce a bijection of the entire vertex sets
before Spoiler chooses the first vertex of the new pair; its mate is
fixed by that bijection. Duplicator wins by preserving equality and
adjacency among all marked pairs through every finite play.
For the joint update above, a winning strategy from the empty position
is equivalent to \(G\equiv_k H\)
(\Cref{lem:wl-game}). The order of these choices is important: the
construction must support a whole bijection, not only a response to an
already selected vertex.

The reduction starts from a fixed exponential-time machine and builds
a pair of polynomial-size graphs. Duplicator will win exactly when the
machine rejects its input. Choosing a decider for the complement of the
source language gives a many-one reduction with the required polarity.
The following three stages explain how an exponential computation can
be represented and tested by a game on these small graphs.

\paragraph{A computation described by its address rules.}
Time and tape position are written in binary, using \(r\) coordinates
in total. We output dimension \(k=r\), so the game has \(r+1\) pair
slots. The circuit gates have a constant number of types and occur
at these addresses. A connection schema lists an allowed set of bits
and an injective map for each coordinate; it represents all connections
whose source addresses lie in the product of those sets.
Carry and borrow cases describe successive times and neighboring tape
positions with polynomially many schemas, without listing the
exponentially many gates.

We next replace the circuit by nonempty sets of one- or two-bit vectors.
Each comparison equates the image of a single scalar affine form at
one endpoint with the image of another form at the other endpoint.
A private two-coordinate buffer makes a true circuit output force the
image \(\{0\}\) at the next input, while a false output can retain the
full image \(\{0,1\}\) and support any nonempty destination image.
Thus the constraints are consistent exactly when every accepting test
is false. Different comparisons may use different vectors from the
same witness set. This individual-image condition, rather than a
simultaneous vector assignment, is the information the later game
will recover. \Cref{sec:computation} proves the circuit and buffer
constructions.

\paragraph{Local affine data with an exact extension property.}
To compress the addressed comparisons, subdivide each schema into a
long wire and arrange the \(r\) address coordinates in a cycle.
A site records a vector on one step of this cycle; a square relates
the two site vectors on opposite sides of one wire link, together
with two scalar values on its other sides. These shared values are
called readouts. Only two consecutive address coordinates occur in
one site or square, so there are polynomially many such blocks and
each has a constant-size space of legal valuations.

The two graph versions will encode a homogeneous and an affine version
of each block equation. An offset translates the homogeneous valuations
of a block to its affine valuations. Offsets on several blocks must
agree on every shared readout. We construct nonempty families of compatible lists on every set of at
most \(r+1\) blocks. Their restrictions are exactly the corresponding
smaller families. Thus every particular list selected by the construction
extends when any one further block is added within this size bound
(\Cref{prop:exact-projection}).

The proof represents offsets using values at the wire types, called
shifts, together with auxiliary values at the corners. A collection of
\(r\) blocks that wraps around the coordinate cycle determines a full
address and constrains a shift through the corresponding witness set.
The original scalar-image equality supplies the preimage needed for
extension. On a component that does not wrap, adjusting the corner
values absorbs a change of shifts while preserving every retained
readout. \Cref{offset:section} makes these corner potentials precise.

\paragraph{Realizing the blocks and reading the game.}
Each graph vertex belongs to exactly one block part. An offset defines
a bijection of that entire part, depending only on the chosen block's
offset. Compatible offsets give an isomorphism on the corresponding
union of parts. A block is active when its part contains a marked vertex.
To announce a global bijection, Duplicator extends the
current list separately for every inactive block, using exact extension.
These choices need not agree with each other: after the announcement,
Spoiler's one selected vertex activates at most one new block.

The converse uses distinguished queries whose responses decode to legal
block valuations. With \(r\) marked sites around one coordinate cycle,
sum the returned vectors to obtain a value at the addressed variable.
The one spare pair moves this ring across the squares of a wire while
retaining each old value until it has been copied. Summing the square
equations cancels the intermediate scalar readouts. Collecting all sums
attained under one winning strategy, and transferring in both directions,
gives exactly the required scalar-image equalities. Responses at different
visits need not coincide.

The general realization makes legal valuations into vertices and uses
private leaves to identify their blocks. The connected subcubic
realization replaces these high-degree parts by valuation paths,
prefix trees and equality matchings for shared readouts; triangles and
private paths identify the vertex sorts. The main additional task is
recognition: a winning strategy must preserve the encoded readouts
even at a history with all slots occupied. Hypothetical continuations
may discard unrelated pairs to test a retained pair, so these tests
establish invariants of the original history without consuming spare
pairs during the ring transfer.

\paragraph{Uniform quantitative bounds.}
For a fixed machine \(M\) running in time \(2^{p(m)}\) on inputs of
length \(m\), the reduction uses
\(r=2(p(m)+2)\) address coordinates and outputs dimension \(k=r\).
The two realizations give the following bounds.
\begin{center}
\begin{tabular}{@{}lll@{}}
\toprule
Realization & Graph order & Adjacency-matrix construction time\\
\midrule
General & \(O_M((r+1)^{12})\) & \(O_{M,p}((r+2)^{80})\)\\
Connected subcubic & \(O_M((r+2)^{16})\) & \(O_{M,p}((r+2)^{160})\)\\
\bottomrule
\end{tabular}
\end{center}
The subscripts indicate dependence only on the fixed machine and
polynomial; the exponents are absolute. Although the subcubic result
already gives qualitative hardness on general graphs, the general
realization has a smaller output and compiler bound. Both bounds are
proved in \Cref{prop:uniform-general,prop:uniform-construction}.
Since the output dimension is polynomial in the source input length,
the two problems remain \(\EXPTIME\)-complete when it is encoded in
unary (\Cref{cor:unary-dimension}).

\paragraph{Organization.}
\Cref{game:section} proves the upper bound and the game characterization.
\Cref{sec:computation} gives the uniform computation and scalar-image
reductions, and \Cref{offset:section} proves exact projection for the
common affine blocks. \Cref{general:section} constructs the general
realization, and \Cref{reduction:section} isolates its reusable properties
and proves the game transfer. \Cref{graph:section} supplies a connected
subcubic realization with those properties. \Cref{uniform:section}
establishes the uniform compiler bounds and completes both reductions.
